\chapter{Java Source File Structure}

Basically, what we refer to as a \bdq{Java Source File} is defined by the \Java programming language as a \bdq{Compilation Unit}\index{compilation unit} \autocite{ORACLE_DOC_LANGUAGE_SPECIFICATION:CompilationUnit}.

\bdq{The \Java Language Specification} distinguishes between \bdq{Ordinary Compilation Units}\index{compilation unit!ordinary}, \bdq{Compact Compilation Units}\index{compilation unit!compact} and \bdq{Modular Compilation Units}\index{compilation unit!modular}. Here, we are only interested in the \bdq{Ordinary Compilation Units}\index{compilation unit!ordinary}.

A \bdq{Modular Compilation Unit}\index{compilation unit!modular} (a \verb#module-info.java#\index{module-info.java} file) defines the module\index{module} structure of an application. The formatting of such a file is described in \tqfullvref{sec:ModuleDefinition} and modules in general are discussed in \tqfullvref{sec:Modularisation}.

\bdq{Compact Compilation Units}\index{compilation unit!compact} or \bdq{compact source files} were introduced with Java~25 and are mainly meant to support writing small scripts in Java, without the overhead of a full-flegded Java program. They should never be used in the context of a bigger Java program or application.

\section{File Encoding}
The \verb#javac# Java compiler expects UTF-8\index{UTF-8} \autocite{UNICODE:CoreSpec,WIKIPEDIA:UTF8} encoded source files. Obviously, ASCII\index{ASCII} \autocite{WIKIPEDIA:ASCII} encoded files will work, too, as well as most ISO-8859\index{ISO-8859} encodings, as long as only the ASCII characters are used.

A resource bundle or properties file\index{resource file}\index{properties file} \autocite{ORACLE_DOC_RESOURCEBUNDLE_CLASS} is expected to be encoded with ISO-8859-1\index{ISO-8859-1} \autocite{WIKIPEDIA:ISO88591}, although this can be changed by setting a system property. Refer to \autocite{ORACLE_DOC_PROPERTYRESOURCEBUNDLE_CLASS} for the details. 

The encodings of other source files depend on their format and usage.

%==============================================================================
% End of file!!
%==============================================================================

\section{Source File Size}
Above a certain number of lines, source code files become unwieldy, but it is difficult to specify a concrete value for the maximum size of a Java source file. First, generated source files tend to have lots of lines because it is often much easier to use boiler-plate code here than to generate methods and call them instead. And in particular when dynamically generated code is used\footnote{Code that is regenerated for each build.}, it would not even hurt the DRY principle: a change to the repeated code can be performed on one single place, the generator.

Second, classes that have lots of attributes usually end up with lots of lines, too: the formatting style and coding guidelines proposed in this document require a least 5~lines for each attribute, 6~lines for each getter method and 11~lines for each setter method, always including all comments and blank lines. This is about 22~lines for each attribute.  One hundred attributes seems to be a lot, but I have seen more than one real life application where database tables had more than that number of columns that each needed to be reflected as an attribute of the class that should reflect the table. I do not think that this is a good design anyway, but if your new Java class has to reflect the legacy data model, it has to have that number of attributes, bringing you to about 2200~lines of code without any methods other than the already mentioned getters and setters.

I found a limit of about 2000~lines to be feasible; if a class becomes larger than this number of lines, you should consider to move some of the code into other classes. This could be a utility class\index{utility class}, a parent class (probably abstract), or the functionality would be better placed with another entity in general. Also externalising some functionality by implementing the command pattern\autocite{Gamma:DesignPatterns} might be an idea. 

\subsubsection{Principles}
\begin{enumerate}[label=P\arabic*.]
\item{A Java source file should not exceed the limit of 2000~lines. As shorter it is, as better.}
\item{Generated source files may have significant more lines.}
\item{Other formatting rules may not be sacrificed only to reduced the number of lines of a source file.}
\end{enumerate}

%==============================================================================
% End of file!!
%==============================================================================

\section{The File Structure}

According to the Java Language Specification\autocite{ORACLE_DOC_LANGUAGE_SPECIFICATION:CompilationUnit} does each (ordinary) compilation unit\index{compilation unit} consist of three parts, each of which is optional:

\begin{itemize}
\item{A package declaration, giving the fully qualified name of the package to which the compilation unit belongs.

A compilation unit that has no package declaration is part of an unnamed package.}
\item{\lstinline|import| declarations that allow classes and interfaces from other packages, and static members of classes and interfaces, to be referred to by using their simple names.}
\item{Top level declarations of classes and interfaces.}
\end{itemize}

Based on that, I define a Java source code file having \textit{five} parts, in the given order:
\begin{enumerate}[nosep]
\item{\nameref{sec:BeginningComment}}
\item{\nameref{sec:PackageStatement}}
\item{\nameref{sec:ImportStatements}}
\item{\nameref{sec:ClassAndInterfaceDeclarations}}
\item{\nameref{sec:ClosingComment}}
\end{enumerate}

\subsection{Beginning Comment}\label{sec:BeginningComment}
A Java project will not consist from Java source files only; aside the already mentioned module definition files(see \tqref{sec:ModuleDefinition}\index{module-info.java}), the package documentation files\index{package-info.java} (refer to \tqfullvref{sec:PackageDocumentation}), there can be various build files, scripts, resource and property files, and even source files other programming languages.

All source files of a project – not only the Java sources – have to start with a comment that gives the copyright\index{copyright} notice and the license\index{license} information\footnote{This is different from the Sun coding conventions\autocite{SUN_CODE_CONVENTIONS:BeginningComments} that recommends to put also the class name and the date into the beginning comment. I omitted the name of the class because this is obvious from the file name in case of a \lstinline|public| class, and it would confuse the reader if the file contains more than one class. The date is considered obsolete.}, if possible at all: it should be obvious that it is sometimes difficult to add this kind of information to a binary resource like an image or a sound file.

For a Java file (a file with a name ending on \verb#*.java#), that beginning comment looks like this:
\begin{lstlisting}[numbers=left,caption={Beginning Comment for a Java file}]
/*
 * ==================================================================
 * Copyright © <Year> <Copyright Notice>
 * ==================================================================
 *
 * <License Notice>
 */
\end{lstlisting}

The same comment block can be used for Groovy\index{Groovy} source files, Kotlin\index{Kotlin} source files, Gradle\index{Gradle} build files (\lstinline|build.gradle|\index{build.gradle}) and generally for all file formats that uses the pattern \lstinline|/* … */| for comments.

A \textit{resource} file\index{resource file}\autocite{ORACLE_DOC_RESOURCEBUNDLE_CLASS} can have XML format or it may contain simple key-value pairs. In the latter form, comments are prepended with the hash symbol (``\verb|#|''), and the beginning comment has the following form: 
\begin{lstlisting}[numbers=left,language=,caption={Beginning Comment for a Resource or a Script file}]
#
# ===================================================================
# Copyright © <Year> <Copyright Notice>
# ===================================================================
#
# <License Notice>
#
\end{lstlisting}

Obviously, this can be used for all other file formats that use the same comment style, like Unix shell scripts\index{shell script}.

For a resource file in XML\index{XML} format\footnote{For an XML file, the beginning comments follows the initial \lstinline|<?xml …?>|, so it is not starting on line~1. That's similar for HTMl files.} (or any other XML file that is part of the sources), as well as for HTML\index{HTML} files, the beginning comment would look like this:
\begin{lstlisting}[numbers=left,firstnumber=5,language=XML,caption={Beginning Comment for XML and HTML files}]
<!--
=====================================================================
 Copyright © <Year> <Copyright Notice>
=====================================================================

 <License Notice>
-->
\end{lstlisting}

\TeX/\LaTeX\index{TeX}\index{LaTeX} uses the per cent sign (``\verb#%#'') to indicate a comment. So the beginning comment for a \TeX/\LaTeX\index{TeX}\index{LaTeX} source file would be like this:
\begin{lstlisting}[numbers=left,language=,caption={Beginning Comment Resource file}]
%
% ===================================================================
% Copyright © <Year> <Copyright Notice>
% ===================================================================
%
% <License Notice>
%
\end{lstlisting}

\verb#<Year>#, \verb#<Copyright Notice># and \verb#<License Notice># has to be replaced by the appropriate texts. The special character \bdq{\copyright} is used not only because is looks better; it tells you immediately whether your environment is capable to deal with UTF-8 encoded files.

The lines with the equals signs (“===”) will end at column~80; refer to chapter \tqfullvref{sec:LineLength}.\footnote{Keep in mind that in this document the line length is set to 70~columns; this means that you cannot just copy and paste the comment blocks into your development environment.}

\subsubsection{Principles}
\begin{enumerate}[label=P\arabic*.]
\item{Each file that is part of the sources for the program/application to build should have a beginning comment.}
\item{That beginning comment must contain the copyright notice with the year and the copytight holder.}
\item{It should contain a license notice (either the license itself or a reference to it).}
\item{No other information should be added to this comment block.}
\end{enumerate}

%==============================================================================
% End of file!!
%==============================================================================

\subsection{Package Statement}\label{sec:PackageStatement}
The first non-comment line of a Java source file is the \lstinline|package| statement\index{package}. According to the Java Language Specification\autocite{ORACLE_DOC_LANGUAGE_SPECIFICATION:UnnamedPackages}, it is optional (resulting in an \bsq{unnamed} or default package), but this is never used for production code. So you must have a \lstinline|package| statement in each and every Java source file.

For the names of packages, see chapter \tqfullvref{sec:Packages}.

\subsubsection{Principles}
\begin{enumerate}[label=P\arabic*.]
\item{Each Java source file has a \lstinline|package| statement.}
\end{enumerate}

%==============================================================================
% End of file!!
%==============================================================================

\subsection{Import Statements}\label{sec:ImportStatements}
Usually several \lstinline|import| statements will follow the \lstinline|package| statement. Only very basic code will not import other stuff than that from the \lstinline|java.lang|\index{java.lang} package. But interfaces with only one method (a so-called \bdq{functional interface}\index{functional interface} \autocite{ORACLE_DOC_FUNCTIONALINTERFACE_ANNOTATION}) may not need any imports at all, same for \bdq{marker interfaces} that does not declare any methods at all. Sometimes you may have also (abstract) parent classes without methods that do not need any \lstinline|import| statements, too.

\keeplisting{8}
But in general, that part of a source file would look like this:
\index{java.lang!String}\index{java.lang!String!format()}\index{java.lang!System}\index{java.lang!System!out}\index{java.awt.peer!CanvasPeer}\index{java.io!InputStream}
\begin{lstlisting}
import static java.lang.String.format;
import static java.lang.System.out;
…

import java.awt.peer.CanvasPeer;
import java.io.InputStream;
…
\end{lstlisting}
Static imports (for methods and constants) have to be placed before the imports for classes. Inside these blocks, the imports should be ordered alphabetically.

\keeplisting{6}
Wildcard imports like
\begin{lstlisting}
import static java.lang.String.*;
…

import java.util.*;
…
\end{lstlisting}
are not allowed and should not even be used for throwaway code. 

Eclipse users can use the function \verb#Source|Organize Imports# from the menu, or the short-key \verb#Shift+Ctrl+O#. This will reorder the import statements, explode the wildcards and remove unused imports; it will even add currently missing imports – given that this is properly configured in the \verb#preferences# at \verb#Java|Code Style|Organize Imports#.

At IntelliJ~IDEA, the function \verb#Code|Optimize Imports# from the menu does something similar; the short key would be \verb#Ctrl+Alt+O#.

\keeplisting{6}
Java~25 introduced a new feature, called \textit{Module Import Declarations} \autocite{JEP511}. It allows to import all public classes and interfaces of a module with a single statement:
\begin{lstlisting}
…
import module java.base;
…
\end{lstlisting}
As it can cause ambiguities, module imports should be avoided.

\subsubsection{Principles}
\begin{enumerate}[label=P\arabic*.]
\item{Place \lstinline|static| imports before non-static ones.}
\item{Sort the \lstinline|import| statements alphabetically.}
\item{Do never use wildcard imports.}
\item{Don't use module imports.}
\item{Remove unused imports. If an import is only required because of a Javadoc reference, change the Javadoc reference to use the fully qualified class name and remove the import.}
\end{enumerate}


%==============================================================================
% End of file!!
%==============================================================================

\subsection{Class and Interface Declarations}\label{sec:ClassAndInterfaceDeclarations}
Each Java source file\footnote{… or \bdq{ordinary compilation unit}\index{compilation unit}} contains at least one top level class or interface declaration; in most cases, it is just one. And it is strongly discouraged that it would be more than one!

Only one of these top level classes or interfaces can be \lstinline|public|, the others have to be package-local (meaning the modifiers \lstinline|private| and \lstinline|protected| are not allowed). If you really have more than one class in a single source file, the \lstinline|public| one should be the first entry, all others will follow, sorted alphabetically by their names.

With the JPMS\index{JPMS|see {Java Platform Module System}}\index{Java Platform Module System} (Java Platform Module System) the need for non-public top-level classes got nearly obsolete and they should be avoided.

If you need classes or interfaces that are closer related to a top-level class, you should consider to define these as \textit{inner} classes or interfaces. Inner classes can be \lstinline|public|, \lstinline|private| or \lstinline|protected|. An inner class is part of the top level class, but when defined as \lstinline|static|, it can be used independently from the enclosing class.

At least since Java~5, there are not only classes and interfaces that are declared in a Java source file. Today, we have

\begin{itemize}[nosep]
\item{\hyperref[subsec:Class]{classes} (keyword \lstinline|class|)}
\item{\hyperref[subsec:Interface]{interfaces} (keyword \lstinline|interface|)}
\item{\hyperref[subsec:Enum]{enums} (keyword \lstinline|enum|)}
\item{\hyperref[subsec:Record]{records} (keyword \lstinline|record|)}
\item{\hyperref[subsec:Annotation]{annotations} (keyword \lstinline|@interface|)}
\end{itemize}

Technically, records and enums are special types of classes, while an annotation\index{annotation} is somehow a special kind of an interface.

Each of these allows different sections in their declarations; the following tables describe the various sections and the order they should appear in. Each section has a header comment, and inside each section, the elements are ordered alphabetically.

\subsubsection{Class}\label{subsec:Class}
\keeplisting{34}
A skeleton for a Java \lstinline|class| may look like this:
\begin{lstlisting}[numbers=left,caption={Class Skeleton}]
public class MyClass 
{
        /*---------------*\
    ====** Inner Classes **==========================================
        \*---------------*/
    …    
        
        /*-----------*\
    ====** Constants **==============================================
        \*-----------*/
    …
        
        /*------------*\
    ====** Attributes **=============================================
        \*------------*/
    …
            
        /*------------------------*\
    ====** Static Initialisations **=================================
        \*------------------------*/
    …
        
        /*--------------*\
    ====** Constructors **===========================================
        \*--------------*/
    …
        
        /*---------*\
    ====** Methods **================================================
        \*---------*/
    …            
}
//  class MyClass
\end{lstlisting}
 
\begin{filecontents}{ClassParts.tbl}
  \begin{longtable}{|l|X|}
  \caption{Parts of a \lstinline|class| declaration} \\
  \hline 
  Part & Description \\ 
  \hline
  \endfirsthead
  \multicolumn{2}{c}%c
  {\tablename\ \thetable\ -- \textit{Continued from previous page}} \\
  \hline 
  Part & Description \\ 
  \hline
  \endhead
  \multicolumn{2}{r}{\textit{Continued on next page}} \\ 
  \endfoot
  \endlastfoot
  Inner Classes (optional) & Given the class defines inner classes, interfaces, enums or records, either \lstinline|private| or \lstinline|public|, they will come first. Internally, they will follow the same rules as for a top-level class. \\ 
  \hline 
  Constants (optional) & This part takes all constants defined by this class (meaning \lstinline|public static final| fields). All values for the constants are known at compile time, otherwise you should place the respective field to the \textit{Static Initialisations} part. \\ 
  \hline 
  Attributes (optional) & Here go the attributes for the class; technically, this part is in fact optional for a class, but usually a class without any attributes does not make much sense. Attributes are \lstinline|private| (under some exceptional circumstances, they can be \lstinline|protected|) and they can be \lstinline|static|, although this should be avoided. \\ 
  \hline 
  Static Initialisations (optional) & If a class declares \lstinline|static| fields that are not mere constants or their initialisation is more complex than just assigning a compile time constant, these fields go into this part.
  
  The fields should be initialised in a \lstinline|static {...}| block, not through an immediate initialiser.
  
  This is also the right place for the \lstinline|serialVersionUID|\index{serialVersionUID} of a seria­lisable\index{java.lang!Serializable} class. \\ 
  \hline 
  Constructors (optional) & All the constructors for the class go into this part. Technically, it is optional if the class will have only an empty \lstinline|public| default constructor, but the recommendation is to code this, too. \\ 
  \hline 
  Methods (optional) & All methods of the class are in this part. Again, this part is technically optional, but classes without methods are barely useful.  \\ 
  \hline 
 \end{longtable} 
\end{filecontents}
\LTXtable{\linewidth}{ClassParts.tbl}

\subsubsection{Interface}\label{subsec:Interface}
\keeplisting{20}
A skeleton for an \lstinline|interface| may look like this:

\begin{lstlisting}[numbers=left,caption={Interface Skeleton}]
public interface MyInterface 
{
        /*---------------*\
    ====** Inner Classes **==========================================
        \*---------------*/
    …
        
        /*-----------*\
    ====** Constants **==============================================
        \*-----------*/
    …
        
        /*---------*\
    ====** Methods **================================================
        \*---------*/
    …
                
}
//  interface MyInterface
\end{lstlisting}
 
\begin{filecontents}{InterfaceParts.tbl}
  \begin{longtable}{|l|X|}
  \caption{Parts of an \lstinline|interface| declaration} \\
  \hline 
  Part & Description \\ 
  \hline
  \endfirsthead
  \multicolumn{2}{c}%
  {\tablename\ \thetable\ -- \textit{Continued from previous page}} \\
  \hline 
  Part & Description \\ 
  \hline
  \endhead
  \multicolumn{2}{r}{\textit{Continued on next page}} \\ 
  \endfoot
  \endlastfoot
  Inner Classes (optional) & Defining inner classes for an interface is discouraged, although there are some valid use cases for this. So it could make sense to define an \lstinline|enum| with flag values used by a method declared by the interface. If the interface defines inner classes or interfaces, they have to be \lstinline|public| – classes have to be \lstinline|static|, too, and they will be the first part. Internally, they will follow the same rules as for top-level classes. \\ 
  \hline 
  Constants (optional) & This part takes all constants defined by this interface (meaning \lstinline|public static final| fields). All values for the constants are known at compile time. \\ 
  \hline 
  Methods (optional) & All methods of the interface are in this part. Again, this part is technically optional, it can be omitted for so-called \textit{marker interfaces}.  \\ 
  \hline 
 \end{longtable} 
\end{filecontents}
\LTXtable{\linewidth}{InterfaceParts.tbl}

\subsubsection{enum}\label{subsec:Enum}
\index{enum}An \lstinline|enum| is a special case of a Java class. Usually it only contains the enumeration values, but it can have more components. A skeleton for a \lstinline|enum| class may look like this:

\begin{lstlisting}[numbers=left,caption={enum Skeleton}]
public enum MyEnum 
{
        /*------------------*\
    ====** Enum Definitions **=======================================
        \*------------------*/
    …
        
        /*-----------*\
    ====** Constants **==============================================
        \*-----------*/
    …
        
        /*------------*\
    ====** Attributes **=============================================
        \*------------*/
    …
        
        /*------------------------*\
    ====** Static Initialisations **=================================
        \*------------------------*/
    …
        
        /*--------------*\
    ====** Constructors **===========================================
        \*--------------*/
    …
        
        /*---------*\
    ====** Methods **================================================
        \*---------*/
    …
                
}
//  enum MyEnum
\end{lstlisting}
 
\begin{filecontents}{EnumParts.tbl}
  \begin{longtable}{|l|X|}
  \caption{Parts of an \lstinline|enum| declaration} \\
  \hline 
  Part & Description \\ 
  \hline
  \endfirsthead
  \multicolumn{2}{c}%
  {\tablename\ \thetable\ -- \textit{Continued from previous page}} \\
  \hline 
  Part & Description \\ 
  \hline
  \endhead
  \multicolumn{2}{r}{\textit{Continued on next page}} \\ 
  \endfoot
  \endlastfoot
  Enum Definitions & This part is mandatory. Each \lstinline|enum| is an instance of the enum class. \\ 
  \hline 
  Constants (optional) & This part takes all additional constants defined by this \lstinline|enum| (meaning \lstinline|public static final| fields). All values for the constants are known at compile time, otherwise you should place the respective field to the \textit{Static Initialisations} part. \\ 
  \hline 
  Attributes (optional) & Here goes the attributes for the \lstinline|enum|, if any. Each attribute has to be \lstinline|private final| and will be initialised by the constructor. \\ 
  \hline 
  Static Initialisations (optional) & If an \lstinline|enum| declares \lstinline|static final| fields with an initialisation that is more complex than just assigning a compile time constant, these fields go into this part.
  
  The fields should be initialised in a \lstinline|static {...}| block. \\
  \hline 
  Constructors (optional) & All the constructors for the enum (rarely there is more than one) go into this part, if any. A constructor for an \lstinline|enum| must be \lstinline|private|. \\ 
  \hline 
  Methods (optional) & All methods of the \lstinline|enum| are in this part.  \\ 
  \hline 
 \end{longtable} 
\end{filecontents}
\LTXtable{\linewidth}{EnumParts.tbl}

Technically, an \lstinline|enum| can also define inner classes, but this is strongly discouraged, and I am not aware of any use case for this.

\subsubsection{Record}\label{subsec:Record}
\index{record}Records had been introduced to Java with version~14 and were finalised with version~16 \autocite{ORACLE_DOC_RECORD,ORACLE_DOC_LANGUAGE_SPECIFICATION:RecordClasses,Eckel:JavaRecords}. Like an \lstinline|enum|, a record is a restricted form of a Java class. It’s ideal for \bdq{plain data carrier} classes that contain data not meant to be altered. Usually it has only the most fundamental methods, such as constructors and accessors.

A skeleton for a Java record class may look like this, with all parts being optional:
\begin{lstlisting}[numbers=left,caption={Record Skeleton}]
public record MyRecord( ... ) 
{
        /*-----------*\
    ====** Constants **==============================================
        \*-----------*/
    …
        
        /*------------*\
    ====** Attributes **=============================================
        \*------------*/
    …
        
        /*------------------------*\
    ====** Static Initialisations **=================================
        \*------------------------*/
    …
        
        /*--------------*\
    ====** Constructors **===========================================
        \*--------------*/
    …
        
        /*---------*\
    ====** Methods **================================================
        \*---------*/
    …
                
}
//  record MyRecord
\end{lstlisting}
 
\begin{filecontents}{RecordParts.tbl}
  \begin{longtable}{|l|X|}
  \caption{Parts of a \lstinline|record| declaration} \\
  \hline 
  Part & Description \\ 
  \hline
  \endfirsthead
  \multicolumn{2}{c}%
  {\tablename\ \thetable\ -- \textit{Continued from previous page}} \\
  \hline 
  Part & Description \\ 
  \hline
  \endhead
  \multicolumn{2}{r}{\textit{Continued on next page}} \\ 
  \endfoot
  \endlastfoot
  Constants (optional) & This part takes all constants defined by this record (meaning \lstinline|public static final| fields). All values for the constants are known at compile time, otherwise you should place the respective field to the \textit{Static Initialisations} part. \\ 
  \hline 
  Attributes (optional) & Here go the attributes for the class; this part is optional, because usually, all attributes are defined as arguments to the \lstinline|record| definition. \\ 
  \hline 
  Static Initialisations (optional) & If a record declares \lstinline|static| fields that are not mere constants or their initialisation is more complex than just assigning a compile time constant, these fields go into this part.
  
  The fields should be initialised in a \lstinline|static {...}| block. \\ 
  \hline 
  Constructors (optional) & If the record needs additional constructors, these go into this part. It is optional, because in most cases no additional constructor is need. \\ 
  \hline 
  Methods (optional) & All additional methods of the record are in this part. \\ 
  \hline 
 \end{longtable} 
\end{filecontents}
\LTXtable{\linewidth}{RecordParts.tbl}

Same as \lstinline|enum|s can records have inner classes, but this is similarly discouraged, and there is no use case for it.

\subsubsection{Annotation}\label{subsec:Annotation}
Annotations had been introduced with Java~5 \autocite{ORACLE_DOC_LANGUAGE_SPECIFICATION:AnnotationInterfaces}. A skeleton for an annotation interface may look like this:

\keeplisting{8}
\begin{lstlisting}[numbers=left,caption={Annotation Skeleton}]
public @interface MyAnnotation 
{
        /*------------*\
    ====** Attributes **=============================================
        \*------------*/
    …
}
//  @interface MyAnnotation
\end{lstlisting}
 
\begin{filecontents}{AnnotationParts.tbl}
  \begin{longtable}{|l|X|}
  \caption{Parts of an annotation declaration} \\
  \hline 
  Part & Description \\ 
  \hline
  \endfirsthead
  \multicolumn{2}{c}%
  {\tablename\ \thetable\ -- \textit{Continued from previous page}} \\
  \hline 
  Part & Description \\ 
  \hline
  \endhead
  \multicolumn{2}{r}{\textit{Continued on next page}} \\ 
  \endfoot
  \endlastfoot
  Attributes (optional) & The attributes for the annotation go here, if any. An annotation without attributes would be a marker annotation. \\ 
  \hline 
 \end{longtable} 
\end{filecontents}
\LTXtable{\linewidth}{AnnotationParts.tbl}

\subsubsection{Principles}
\begin{enumerate}[label=P\arabic*.]
\item{Only one public top-level class (interface, record, enum, …) per compilation unit. Avoid non-public top-level classes.}
\item{Group the parts of the class as shown above. Remove the comments for unused sections.}
\end{enumerate}


%==============================================================================
% End of file!!
%==============================================================================


%==============================================================================
% End of file!!
%==============================================================================


%==============================================================================
% End of file!!
%==============================================================================
